\section{Mesh generation}\label{sec:mesh-generation}

A Zeffiro forward model starts from closed tissue surfaces and fills the
interior with tetrahedra. Each tetrahedron is labelled by a compartment
(scalp, skull, CSF, grey matter, \ldots) so later stages can assign
conductivity and decide which tetrahedra may contain neural sources.

This chapter describes the live pipeline in \ccode{src/mesh}. The Mesh tool
window title is \emph{ZEFFIRO Interface: Mesh tool}. The button
\textbf{Create FEM mesh} does not call \ccode{zef\_create\_fem\_mesh} alone;
it runs the wrapper \ccode{zef\_create\_finite\_element\_mesh}
(optional surface downsample, \ccode{zef\_process\_meshes}, volume fill,
\ccode{zef\_postprocess\_fem\_mesh}). There is no live \textbf{make\_all}
button.

Scripted equivalent:
\begin{lstlisting}[style=unnumberedlines, language=Matlab]
zef = zef_create_finite_element_mesh(zef);
% or, without the downsample wrapper:
zef = zef_process_meshes(zef);
zef = zef_create_fem_mesh(zef);
zef = zef_postprocess_fem_mesh(zef);
zef = zef_update(zef);
\end{lstlisting}

Surfaces must already exist as \ccode{zef.<tag>\_points} /
\ccode{zef.<tag>\_triangles} (Import segmentation, or a \ccode{.zef} import
script). Length units of surfaces, sensors, and \ccode{mesh\_resolution}
must match; the lattice builder does not convert millimetres to metres.

\subsection{Pipeline}

\begin{enumerate}
\item \ccode{zef\_process\_meshes} gathers active compartments into
      \ccode{zef.reuna\_p} / \ccode{reuna\_t}, applies scaling, Euler
      rotations (degrees about the pre-affine centroid), and affine
      transforms, and prepares sensors.
\item \ccode{zef\_create\_fem\_mesh} fills the bounding box with a cubic
      lattice split into tetrahedra, then labels them with solid-angle
      inside tests (\ccode{src/compartments}).
\item Optional surface / volume / adaptive refinement, with relabeling when
      \ccode{zef.mesh\_relabeling} is on.
\item \ccode{zef\_postprocess\_fem\_mesh} optionally smooths, flips inverted
      tetrahedra (\ccode{zef\_tetra\_turn}), writes \ccode{zef.sigma} from
      the parameter profile, and sets \ccode{brain\_ind}.
\end{enumerate}

\textbf{Lattice.} Default \ccode{initial\_mesh\_mode} is \(1\): five
tetrahedra per cube, with a parity-dependent stencil so shared faces use the
same diagonal. Mode \(2\) uses six tetrahedra and one stencil for every cube.
A compartment with \ccode{<tag>\_sources == -1} is a perfectly matched layer
(PML); the lattice then comes from \ccode{zef\_pml\_mesh}.

\textbf{Labeling.} \ccode{mesh\_labeling\_approach} \(1\) stores the eight
cube-corner indices per tetrahedron; approach \(2\) stores the four tet
vertices. Exterior tetrahedra (vertices not all inside some compartment) are
discarded. Overlaps are resolved with \ccode{zef.priority\_mode} and
per-compartment labeling priority.

\subsection{Fields of \code{zef} needed by mesh generation}%
\label{ssec:meshing-fields}

Set fields with
\begin{lstlisting}[style=unnumberedlines, language=Matlab]
zef.field_name = field_value;
\end{lstlisting}
Defaults below are from \ccode{src/app/zef\_init.m}. MATLAB does not
enforce the listed sizes at runtime.

\zeffield{compartment\_tags}{1,:}{cell}{}{%
	Short tags (for example \ccode{d1}, \ccode{w}) used to build dynamic
	field names. The segmentation-tool table rows are stored in reverse
	order relative to this cell.
}

\zeffield{\comptag\_on}{1,1}{logical or double}{}{%
	Whether the compartment participates in meshing and plotting. Off
	compartments are dropped by \ccode{zef\_update}.
}

\zeffield{\comptag\_scaling}{1,1}{double}{}{%
	Uniform scale applied in \ccode{zef\_process\_meshes}.
}

\zeffield{\comptag\_x\_correction}{1,1}{double}{}{%
	Translation in \(x\) (same unit as the surface points).
}

\zeffield{\comptag\_y\_correction}{1,1}{double}{}{%
	Translation in \(y\).
}

\zeffield{\comptag\_z\_correction}{1,1}{double}{}{%
	Translation in \(z\).
}

\zeffield{\comptag\_xy\_rotation}{1,1}{double}{}{%
	Rotation in degrees about the compartment centroid, \(xy\) plane,
	applied before the affine cell.
}

\zeffield{\comptag\_yz\_rotation}{1,1}{double}{}{%
	Rotation in degrees, \(yz\) plane.
}

\zeffield{\comptag\_zx\_rotation}{1,1}{double}{}{%
	Rotation in degrees, \(zx\) plane.
}

\zeffield{\comptag\_points}{:,3}{double}{}{%
	Surface vertices of this compartment (typically millimetres).
}

\zeffield{\comptag\_triangles}{:,3}{double}{}{%
	Surface triangles (1-based indices into \ccode{\comptag\_points}).
}

\zeffield{\comptag\_submesh\_ind}{1,:}{double}{}{%
	Cumulative last-face indices of concatenated submeshes (for example
	left/right cortex). Empty means a single patch.
}

\zeffield{\comptag\_points\_inf}{:,3}{double}{}{%
	Inflated copy of the surface used to keep sources slightly inside
	tissue. Filled by \ccode{zef\_downsample\_surfaces} /
	\ccode{zef\_inflate\_surface} unless \ccode{zef.bypass\_inflate} is
	true.
}

\zeffield{\comptag\_sources}{1,1}{double}{}{%
	Source / activity code used when building \ccode{zef.reuna\_type}:
	\(-1\) PML / bounding box, \(0\) inactive, \(1\) constrained field,
	\(2\) unconstrained field, \(3\) active surface. Lead-field
	``brain'' tetrahedra are those with codes \(1\) or \(2\).
}

\zeffield{\comptag\_sigma}{1,1}{double}{}{%
	Isotropic conductivity in siemens per metre, copied onto tetrahedra
	during postprocess.
}

\zeffield{\comptag\_affine\_transform}{1,:}{cell}{}{%
	Sequence of \(4\times 4\) affines applied as homogeneous rows times
	\(A^\top\) in \ccode{zef\_process\_meshes}.
}

\zeffield{mesh\_resolution}{1,1}{double}{mustBePositive}{%
	Cube edge length of the initial lattice, in the same unit as the
	surfaces. Default \(3\). Smaller values produce more tetrahedra.
}

\zeffield{initial\_mesh\_mode}{1,1}{double}{}{%
	\(1\): five tetrahedra per cube (parity stencils). \(2\): six
	tetrahedra per cube. Default \(1\).
}

\zeffield{downsample\_surfaces}{1,1}{logical}{}{%
	Mesh tool \textbf{Resample surf.} checkbox. When true, the create
	wrapper calls \ccode{zef\_downsample\_surfaces} first. Default on
	in \ccode{zef\_init}.
}

\zeffield{max\_surface\_face\_count}{1,1}{double}{}{%
	Target face count scale for downsampling (times relative resolution).
}

\zeffield{mesh\_smoothing\_on}{1,1}{logical}{}{%
	If true, postprocess runs \ccode{zef\_smoothing\_step} (Taubin).
	Default off.
}

\zeffield{mesh\_smoothing\_repetitions}{1,1}{double}{}{%
	Number of smoothing sweeps. Default \(1\).
}

\zeffield{mesh\_optimization\_repetitions}{1,1}{double}{}{%
	Maximum flips in \ccode{zef\_tetra\_turn} while inverted or
	poor-quality tetrahedra remain. Default \(10\).
}

\zeffield{mesh\_optimization\_parameter}{1,1}{double}{}{%
	Quality threshold passed to \ccode{zef\_tetra\_turn} (default
	\(10^{-5}\)).
}

\zeffield{mesh\_labeling\_approach}{1,1}{double}{}{%
	\(1\): label from the eight cube corners. \(2\): label from the four
	tet vertices. Default \(1\). Edited in Forward and inverse
	processing options, not on the Mesh tool itself.
}

\zeffield{mesh\_relabeling}{1,1}{logical}{}{%
	Re-run labeling after each refinement or smoothing pass. Default on.
}

\zeffield{refinement\_on}{1,1}{logical}{}{%
	Master switch for the Mesh tool \textbf{Refinement} checkbox.
}

\zeffield{refinement\_surface\_on}{1,1}{logical}{}{%
	Split tetrahedra that meet selected compartment surfaces
	(\ccode{zef\_refinement\_step}).
}

\zeffield{refinement\_surface\_compartments}{1,:}{double}{}{%
	Which compartments to refine. \(-1\) means all source compartments
	(\ccode{\_sources} in \(\{1,2\}\)).
}

\zeffield{refinement\_volume\_on}{1,1}{logical}{}{%
	Uniform 4-to-1 (actually 8-tet) volume split of selected domains
	(\ccode{zef\_mesh\_refinement}).
}

\zeffield{pml\_outer\_radius}{1,1}{double}{}{%
	PML cube half-width. Unit \(1\): relative to \(\max|\mathrm{reuna}_p|\).
	Unit \(2\): absolute. Used when some \ccode{\_sources == -1}.
}

After a successful create/postprocess you should have \ccode{zef.nodes}
(\(V\times 3\)), \ccode{zef.tetra} (\(T\times 4\), 1-based),
\ccode{zef.domain\_labels} (\(T\times 1\)), and typically \ccode{zef.sigma}.
The lead-field chapter assumes those fields exist.
